\begin{python0}
from solutions import *; clear()
\end{python0}
\sectionthree{Operators <op> and <op>=}

Defining \textit{+} does not give you \textit{+=} automatically.

Here's some advice: define + in terms of \textit{+=},
\textit{-} in terms of \textit{-=}, etc. For instance:

\begin{consolethree}[escapeinside=||]
class Int
{
public:
    ...
    Int & operator+=(const Int & c)
    {
        x_ += c.x_;
        return (*this);
    }
    Int operator+(Int c) const
    {
        return (Int(*this) += c);
    }
    ...
private:
    int x_;
};
\end{consolethree}

A \textbf{nonmember function} is just a function not in a class (i.e.,
not a method). If the function works with objects from a particular
class, then the function is included with the class. You should

\begin{enumerate}
\item Put the function prototype in the header file (\textit{*.h})
\item Define the function in the implementation file (\textit{*.cpp})
\item Note that the function prototype is outside the class declaration
\end{enumerate}

Since the function prototype is out the class, it does \textbf{not} have
access to the private members of the class. You can overwrite this for
efficiency (later). Note that here \textit{+=} is more efficient than \textit{+}, same
for integer \textit{+} and \textit{+=}. Also note that we want to do

\textit{(a += b) += b;}

\textit{c = (a += b); // i.e., a+=b; c=a;}

That's why \textit{operator+=} returns a reference to
\textit{*this}.

\begin{consolethree}[escapeinside=||]
class Int
{
public:
    ...
    |{\bf Int \&}| operator+=(const C & c)
    {
        x += c.x;
        |{\bf return (*this);}|
    }
    ...
};
\end{consolethree}

And \textit{operator+} can be written like this:

\begin{consolethree}[escapeinside=||]
class Int
{
public:
    ...
    Int operator+(const Int c) const
    {
        |{\bf return (Int(*this) += c);}|
    }
    ...
};
\end{consolethree}


